\documentclass[11pt]{amsart}

\usepackage[T1]{fontenc}
\usepackage[utf8]{inputenc}
\usepackage{amsmath,amssymb,amsthm}
\usepackage[margin=1in]{geometry}
\usepackage{hyperref}

\newcommand{\OO}{\mathbb{O}}
\newcommand{\RR}{\mathbb{R}}
\newcommand{\ZZ}{\mathbb{Z}}
\newcommand{\Nrm}{N}

\title[Historical appendix: integral Cayley numbers 1923--1946]%
{Draft appendix for\\
\emph{An order on the Leech lattice from a $\mathbb{Z}_3$-symmetric
triple-octonion product}\\
\large Historical note: integral Cayley numbers, 1923--1946}
\author{Claude Opus 4.7}
\address{Anthropic}
\thanks{Drafted by the AI agent Claude Opus 4.7 (Anthropic) at the
direction of Jens K\"oplinger, as part of the systematic investigation
recorded in this repository (\texttt{prompt\_logs/}).  The version that
appears in the main paper is owned by the principal investigator.}
\date{23 May 2026}

\begin{document}
\maketitle

% -----------------------------------------------------------------
% Body intended for v4 as an appendix.  Target length in the
% main paper: <= 2 pages (target 1 page).
% -----------------------------------------------------------------

\section*{Historical note: integral Cayley numbers, 1923--1946}
\label{app:history}

The construction this paper builds on -- an integral-octonion maximal
order supported on the $E_8$ lattice, repaired by a transposition of two
imaginary basis units -- was developed across four primary papers between
1923 and 1946.  Each was obtained in the original and re-derived in exact
arithmetic in this project (\texttt{python\_project/src/verify\_\{dickson%
\_1923,kirmse\_1924,mahler\_1942,coxeter\_1946\}.py}; verification notes in
\texttt{paper/reviews/}).

\noindent\textbf{Dickson (1923) \cite{Dickson1923}} gave the first
construction: the Cayley--Dickson doubled product and an explicit
8-element basis of an order $O_D$, with all 64 basis products closing.
His Theorem~XV states three maximal orders; the true count is
\emph{seven}.  The cause is in his derivation (p.~321): he assumes the
orders contain the Hurwitz quaternions, and exactly three of the seven
do.

\noindent\textbf{Kirmse (1924) \cite{Kirmse1924}} took up the problem
independently.  His multiplication table is itself a correct octonion
algebra, and he derived a correct \emph{necessary} form for the
maximal orders -- thirty $E_8$-type lattices realise it -- without
proving sufficiency.  Among the eight of those lattices that are
stable under multiplication by units (an unstated narrowing of the
criterion, in our reconstruction), he exhibited the one, $J_1$
(p.~70), that is not an order.  The remaining seven of his eight are
the genuine maximal orders.

\noindent\textbf{Mahler (1942) \cite{Mahler1942}} developed the
arithmetic of a maximal order: for every Cayley number~$X$ there is an
integral~$G$ with $\Nrm(X-G)\le\tfrac12$ (sharper than Dickson's $5/4$),
and every left or right ideal is principal, generated by an integral
Cayley number of norm $1$, $2$, or $4$.

\noindent\textbf{Coxeter (1946) \cite{Coxeter1946}} identifies Kirmse's
error with the same non-closure witness and records that the remedy is
due to R.\,H.~Bruck: ``Bruck's domain $J$ can be derived from $J_1$ by
transposing two of the $i$'s'' (\S4).  This is a transposition of two
imaginary basis units, a linear involution of~$\RR^8$; the
$\binom72 = 21$ such transpositions yield the seven maximal orders in
seven sets of three.  His \S\S6--13 develop the $E_8/5_{21}$ geometry of
the corrected order.  A postscript (\S14) records his after-the-fact
awareness of Dickson and Mahler and identifies Dickson's undercount with
exactly the cause above.

The count converges $3 \to 8 \to 7$.  The \emph{Coxeter--Bruck
transposition} -- a transposition of two imaginary coordinate axes -- is
the historical precedent for the present paper's~$\sigma$.  The two are
the same kind of map, a linear involution of~$\RR^8$, and differ only in
what they are asked about: the Coxeter--Bruck transposition acts on a
maximal-order candidate; $\sigma$ acts on the sublattice~$Ls$ of Wilson's
Leech construction.

% -----------------------------------------------------------------
\begin{thebibliography}{4}
% -----------------------------------------------------------------

\bibitem[Dickson1923]{Dickson1923}
L.\,E.~Dickson, \emph{A new simple theory of hypercomplex integers},
J.~Math.\ Pures Appl.\ (9) \textbf{2} (1923), 281--326.

\bibitem[Kirmse1924]{Kirmse1924}
J.~Kirmse, \emph{\"Uber die Darstellbarkeit nat\"urlicher ganzer Zahlen
als Summen von acht Quadraten \ldots}, Ber.\ Verh.\ S\"achs.\ Akad.\
Wiss.\ Leipzig, Math.-Phys.\ Kl.\ \textbf{76} (1924), 63--82.

\bibitem[Mahler1942]{Mahler1942}
K.~Mahler, \emph{On ideals in the Cayley--Dickson algebra},
Proc.\ Roy.\ Irish Acad.\ Sect.\ A \textbf{48} (1942), 123--133.

\bibitem[Coxeter1946]{Coxeter1946}
H.\,S.\,M.~Coxeter, \emph{Integral Cayley numbers},
Duke Math.\ J.\ \textbf{13} (1946), 561--578.

\end{thebibliography}

\end{document}
